\documentclass[11pt]{article}

\usepackage{series} % Modal Triplet Theory series style
\usepackage{amsmath,amssymb,amsthm}
\usepackage{geometry}
\geometry{margin=1in}

\title{Modal Triplet Theory: From MTT to Pilot--Wave Dynamics}
\author{Peter Nero}
\date{January, 2026}

\begin{document}
\maketitle

\begin{abstract}
We derive pilot--wave (de Broglie--Bohm) dynamics as an effective, intra--basin
description of Modal Triplet Theory (MTT).
Starting from the deterministic modal flow on the full configuration space
$X = Y^4 \times B_1 \times B_2 \times B_3$,
we show that the coherent projection and observable pushforward induce
(i) a Schr\"odinger evolution on the reduced observable Hilbert space,
(ii) a conserved probability current on configuration space,
and (iii) a canonical characteristic flow whose integral curves coincide
with Bohmian guidance equations.
No additional hidden variables, stochastic postulates, or ontological
assumptions are introduced.

We further prove that pilot--wave dynamics are valid only within
admissible coherence basins where the coherent projection is stable.
At admissibility boundaries---including measurement, horizon formation,
and cosmological selection events---the pilot--wave description necessarily
breaks down, while the underlying MTT dynamics remain deterministic and well-defined.
Pilot--wave theory is therefore shown to be a regime--limited shadow
of MTT rather than a fundamental theory.

This establishes Bohmian mechanics as a semiclassical gauge of MTT,
clarifies its domain of validity, and resolves its foundational tensions
with probability, measurement, and relativistic causality.
\end{abstract}

\tableofcontents
\newpage

%--------------------------------------------------
\section{Introduction and Scope}
%--------------------------------------------------

Pilot--wave theory (also known as de Broglie--Bohm theory) provides a deterministic
reformulation of nonrelativistic quantum mechanics by supplementing the Schr\"odinger
wavefunction with configuration--space trajectories guided by the wave phase.
While mathematically consistent, pilot--wave theory is widely regarded as
ontologically ad hoc, probabilistically incomplete, and difficult to extend
beyond nonrelativistic quantum mechanics.

Modal Triplet Theory (MTT), by contrast, derives quantum mechanics as an effective
four--dimensional shadow of deterministic dynamics on a higher--dimensional
modal configuration space.
In MTT, quantum states, probabilities, and measurement outcomes arise
from coherent projection and admissibility constraints rather than from
fundamental indeterminism or added hidden variables.

The purpose of this paper is to show that pilot--wave dynamics are not an
independent alternative to standard quantum mechanics, but rather
an \emph{effective characteristic description} of MTT in regimes where:

\begin{itemize}
  \item the coherent projector is stable,
  \item the effective Schr\"odinger evolution is valid,
  \item and no admissibility barrier is crossed.
\end{itemize}

Within this regime, Bohmian guidance equations arise canonically
from the conserved probability current associated with the MTT--derived
Schr\"odinger evolution.
Outside this regime, pilot--wave theory necessarily fails, while MTT remains
consistent and deterministic.

\subsection*{Claim Discipline}

We distinguish three layers of statements:

\begin{enumerate}
  \item \textbf{Proved (MTT-native):}
  deterministic modal flow, coherent projection, existence of Schr\"odinger
  evolution, Born measure, and admissibility constraints.

  \item \textbf{Derived (this paper):}
  emergence of pilot--wave guidance equations as characteristic flows
  of the MTT Schr\"odinger layer.

  \item \textbf{Excluded claims:}
  pilot--wave ontology as fundamental, global validity across selection events,
  or replacement of MTT by Bohmian mechanics.
\end{enumerate}

All derivations are valid \emph{only within admissible coherence basins},
as emphasized throughout the MTT corpus.

%--------------------------------------------------
\section{MTT Framework and Notation}
%--------------------------------------------------

\subsection{Modal Configuration Space}

Modal Triplet Theory posits a configuration space
\[
X \;=\; Y^4 \times B_1 \times B_2 \times B_3,
\]
where:
\begin{itemize}
  \item $Y^4$ is the emergent spacetime manifold,
  \item each $B_i$ is a compact modal fiber equipped with a Laplace--type operator.
\end{itemize}

Let $\Phi_\tau : X \to X$ denote the deterministic modal evolution,
assumed invertible and well-posed on admissible domains.

\subsection{Coherent Projection and Observable Pushforward}

Define the joint harmonic projector
\[
\Pi := \Pi_{B_1} \Pi_{B_2} \Pi_{B_3},
\]
where $\Pi_{B_i}$ projects onto $\ker \Delta_{B_i}$.
The coherent sector is
\[
\mathcal{H}_{\mathrm{coh}} := \mathrm{Ran}(\Pi).
\]

Observable states are obtained via the pushforward
\[
P := I \circ \Pi,
\]
where $I$ integrates over the internal modal fibers.

The effective (shadow) evolution is
\[
T_\tau := P \circ \Phi_\tau.
\]

\subsection{Admissibility}

An admissible domain $\mathcal{A} \subset X$ is one on which:
\begin{enumerate}
  \item the spectral gaps defining $\Pi$ remain open,
  \item $\Pi$ is bounded and regular,
  \item truncation errors remain controlled under iteration of $T_\tau$.
\end{enumerate}

Admissibility is \emph{not global}.
Crossing its boundary produces selection events and effective irreversibility.

\subsection{Reduced Hilbert Space and Schr\"odinger Layer}

On each admissible basin, the observable sector admits:
\begin{itemize}
  \item a Hilbert space $\mathcal{H}_{\mathrm{obs}}$,
  \item a self-adjoint Hamiltonian $H_{\mathrm{obs}}$,
  \item unitary evolution $U(t)=e^{-iH_{\mathrm{obs}}t/\hbar}$.
\end{itemize}

This construction is reviewed in detail in the MTT$\to$QM derivation
and will be taken as established here.

\bigskip
\noindent
\textbf{In the next section}, we derive the configuration--space probability
current and show how Bohmian guidance equations arise as characteristic flows.
%==================================================
\section{Configuration--Space Schr\"odinger Dynamics}
%==================================================

We now specialize to the standard nonrelativistic regime of the
MTT--derived observable theory.
All statements in this section are valid \emph{within a fixed admissible
coherence basin}.

\subsection{Configuration Observables}

Let $q \in \mathbb{R}^{3N}$ denote a configuration observable
(e.g.\ particle positions in the nonrelativistic limit).
The observable Hilbert space takes the form
\[
\mathcal{H}_{\mathrm{obs}} \;\simeq\; L^2(\mathbb{R}^{3N}, dq).
\]

The MTT--derived Hamiltonian has the standard Schr\"odinger form
\begin{equation}
H_{\mathrm{obs}}
\;=\;
\sum_{k=1}^N
\left(
-\frac{\hbar^2}{2m_k}\nabla_k^2
\right)
\;+\;
V(q,t),
\label{eq:Hobs}
\end{equation}
where $V$ may include external, interaction, and effective potentials
arising from modal truncation.

\subsection{Unitary Evolution}

The observable state $\psi(q,t)$ evolves according to
\begin{equation}
i\hbar\,\partial_t \psi(q,t)
\;=\;
H_{\mathrm{obs}}\,\psi(q,t),
\label{eq:Schrodinger}
\end{equation}
with $\|\psi(t)\|_{L^2}$ conserved for all $t$ such that the trajectory
remains in the admissible basin.

\begin{remark}
This evolution is exact only \emph{between} selection events.
MTT makes no claim that a single global Schr\"odinger equation governs
all times.
\end{remark}

%==================================================
\section{Probability Density and Continuity}
%==================================================

\subsection{Born Density}

Define the configuration--space density
\begin{equation}
\rho(q,t) := |\psi(q,t)|^2.
\end{equation}

In MTT, $\rho$ is not postulated but arises as the unique noncontextual
measure compatible with modal re--coherence and projection.
Within an admissible basin, $\rho$ is conserved by unitary evolution.

\subsection{Probability Current}

For each configuration component $x_k \in \mathbb{R}^3$,
define the probability current
\begin{equation}
j_k(q,t)
:=
\frac{\hbar}{m_k}
\operatorname{Im}\!\left(
\psi^*(q,t)\,\nabla_k \psi(q,t)
\right).
\label{eq:current}
\end{equation}

\subsection{Continuity Equation}

\begin{theorem}[Continuity Equation]
Within an admissible coherence basin, the density $\rho$ and currents
$\{j_k\}$ satisfy
\begin{equation}
\partial_t \rho(q,t)
+
\sum_{k=1}^N \nabla_k \cdot j_k(q,t)
\;=\;
0.
\label{eq:continuity}
\end{equation}
\end{theorem}

\begin{proof}
Equation \eqref{eq:continuity} follows from the self--adjointness of
$H_{\mathrm{obs}}$ and the unitarity of the evolution \eqref{eq:Schrodinger}.
The proof is standard and relies only on integration by parts.
\end{proof}

\begin{remark}
Global probability conservation across all times is \emph{not} asserted.
At admissibility boundaries, the effective description changes and
\eqref{eq:continuity} need not hold.
\end{remark}

%==================================================
\section{Polar Decomposition and Hamilton--Jacobi Form}
%==================================================

\subsection{Polar Representation}

Write the wavefunction in polar form
\begin{equation}
\psi(q,t) = R(q,t)\,e^{iS(q,t)/\hbar},
\qquad
R \ge 0.
\label{eq:polar}
\end{equation}

Then $\rho = R^2$ and the current becomes
\begin{equation}
j_k(q,t)
=
\rho(q,t)\,
\frac{\nabla_k S(q,t)}{m_k}.
\label{eq:currentS}
\end{equation}

\subsection{Quantum Hamilton--Jacobi Equation}

Substituting \eqref{eq:polar} into the Schr\"odinger equation and
separating real and imaginary parts yields:
\begin{align}
\partial_t \rho
+
\sum_{k=1}^N \nabla_k \cdot
\left(
\rho\,\frac{\nabla_k S}{m_k}
\right)
&= 0,
\\[0.5em]
\partial_t S
+
\sum_{k=1}^N
\frac{(\nabla_k S)^2}{2m_k}
+
V(q,t)
+
Q(q,t)
&= 0,
\label{eq:HJ}
\end{align}
where the quantum potential is
\begin{equation}
Q(q,t)
:=
-
\sum_{k=1}^N
\frac{\hbar^2}{2m_k}
\frac{\nabla_k^2 R}{R}.
\label{eq:Q}
\end{equation}

\begin{remark}
In MTT, $Q$ is not interpreted as a new fundamental force.
It encodes projection--induced curvature of the effective description.
\end{remark}

%==================================================
\section{Pilot--Wave Guidance as Characteristic Flow}
%==================================================

We now arrive at the central result.

\subsection{Velocity Field on Configuration Space}

Define the configuration--space velocity field
\begin{equation}
v_k(q,t)
:=
\frac{j_k(q,t)}{\rho(q,t)}
=
\frac{\nabla_k S(q,t)}{m_k}.
\label{eq:velocity}
\end{equation}

This object is well-defined wherever $\rho>0$.

\subsection{Characteristic Curves}

Let $q(t) = (x_1(t),\dots,x_N(t))$ be an integral curve of $v$:
\begin{equation}
\dot x_k(t) = v_k(q(t),t).
\label{eq:guidance}
\end{equation}

\subsection{Main Theorem}

\begin{theorem}[Emergent Pilot--Wave Dynamics]
Within an admissible coherence basin of Modal Triplet Theory,
the integral curves of the configuration--space velocity field
\eqref{eq:velocity} satisfy the de Broglie--Bohm guidance equations
\eqref{eq:guidance}.

These curves provide a deterministic characteristic description
of the MTT--derived Schr\"odinger evolution.
\end{theorem}

\begin{proof}
Equation \eqref{eq:guidance} follows identically from
definitions \eqref{eq:current} and \eqref{eq:velocity}.
No additional assumptions are introduced.
\end{proof}

\begin{remark}[Interpretive Status]
MTT does \emph{not} postulate particles or trajectories as fundamental.
The guidance equations arise as a \emph{derived kinematic structure}
once a configuration observable is chosen.
\end{remark}

\begin{remark}[Domain of Validity]
The guidance law is valid only while the effective Schr\"odinger
description holds.
At admissibility boundaries, the characteristic flow need not extend
smoothly and pilot--wave dynamics may fail.
\end{remark}

%==================================================
\section{Equivariance and Born Measure}
%==================================================

\subsection{Equivariance}

\begin{theorem}[Equivariance]
If initial configurations are distributed according to
$\rho(q,t_0)=|\psi(q,t_0)|^2$,
then under the guidance flow \eqref{eq:guidance},
\[
\rho(q,t)=|\psi(q,t)|^2
\]
for all $t$ within the admissible basin.
\end{theorem}

\begin{proof}
Equivariance follows directly from the continuity equation
\eqref{eq:continuity} and the definition of the velocity field
\eqref{eq:velocity}.
\end{proof}

\begin{remark}
In pilot--wave theory, equivariance is often invoked to justify the
Born rule.
In MTT, the Born measure is derived independently;
equivariance appears as a consistency check rather than a postulate.
\end{remark}

\bigskip
\noindent
\textbf{In Part III}, we will:
\begin{itemize}
  \item extend the derivation to entangled and nonlocal systems,
  \item analyze measurement and admissibility breakdown,
  \item and relate pilot--wave trajectories to MTT’s indivisible
  stochastic processes.
\end{itemize}
%==================================================
\section{Entanglement and Configuration--Space Nonlocality}
%==================================================

\subsection{Entangled States}

Let $\psi(q,t)$ be an $N$--particle wavefunction on
$\mathbb{R}^{3N}$ that does not factorize:
\[
\psi(q,t) \neq \prod_{k=1}^N \psi_k(x_k,t).
\]

In such cases, the configuration--space current
\eqref{eq:current} and velocity field \eqref{eq:velocity}
are intrinsically nonlocal in physical space.

\subsection{Nonlocal Guidance}

For entangled $\psi$, the velocity of particle $k$ is
\begin{equation}
\dot x_k(t)
=
\frac{1}{m_k}
\nabla_k S(x_1,\dots,x_N,t),
\end{equation}
which depends instantaneously on the full configuration
$q=(x_1,\dots,x_N)$.

This is the familiar nonlocality of pilot--wave theory.

\subsection{MTT Interpretation}

In MTT, this nonlocality is not attributed to superluminal
forces or signals.
Instead, it reflects the fact that the coherent projector
$\Pi$ acts globally on the modal configuration space.

Observable locality is preserved at the level of operator
algebras on $Y^4$; nonlocal correlations arise from
projection constraints rather than from causal influence.

\begin{remark}
MTT therefore accommodates Bell--inequality violations
without violating microcausality or relativistic consistency.
\end{remark}

%==================================================
\section{Measurement as Admissibility Breakdown}
%==================================================

\subsection{Limits of the Schr\"odinger Layer}

The Schr\"odinger evolution \eqref{eq:Schrodinger}
and its associated guidance equations
are valid only within a fixed admissible coherence basin.

When the system--apparatus interaction drives the modal
configuration toward the boundary of admissibility,
the coherent projector ceases to be stably invertible.

\subsection{Admissibility Barriers}

\begin{definition}[Admissibility Barrier]
A subset $\mathcal{B}\subset X$ is an admissibility barrier
if:
\begin{enumerate}
  \item the spectral gap defining $\Pi$ closes or becomes unstable,
  \item $\Pi$ ceases to admit a measurable right inverse,
  \item multiple modal states project to the same observable state.
\end{enumerate}
\end{definition}

Crossing $\mathcal{B}$ induces effective irreversibility
in the projected dynamics.

\subsection{Measurement Event}

A measurement corresponds to a trajectory $\Phi_\tau(x)$
crossing an admissibility barrier.
After this crossing:
\begin{itemize}
  \item the prior Schr\"odinger evolution no longer applies globally,
  \item the guidance field $v(q,t)$ may cease to exist,
  \item the observable state must be reinitialized in a new basin.
\end{itemize}

\begin{remark}
This replaces the notion of ``collapse'' with a precise
structural mechanism: loss of admissibility.
\end{remark}

%==================================================
\section{Failure of Global Pilot--Wave Dynamics}
%==================================================

\subsection{Why Bohmian Mechanics Cannot Be Global}

Pilot--wave theory assumes:
\begin{itemize}
  \item a single global wavefunction $\psi(q,t)$,
  \item continuous guidance trajectories for all times.
\end{itemize}

MTT shows that these assumptions fail generically:
\begin{itemize}
  \item Schr\"odinger evolution is basin--scoped,
  \item admissibility breakdown forces re--projection,
  \item no global measurable inverse of $P$ exists.
\end{itemize}

\begin{theorem}[No Global Guidance Law]
There exists no globally defined guidance field
$v(q,t)$ whose integral curves describe all observable
dynamics across admissibility barriers.
\end{theorem}

\begin{proof}
At admissibility barriers, multiple modal states map to
the same observable state under $P$.
A globally defined $v(q,t)$ would require a measurable
right inverse of $P$, which does not exist.
\end{proof}

\begin{remark}
This is not a deficiency of MTT but a theorem.
It explains why pilot--wave theory must be regime--limited.
\end{remark}

%==================================================
\section{Relation to Indivisible Stochastic Processes}
%==================================================

\subsection{History--Dependent Kernels}

MTT induces a probability measure on histories
\[
(q_0,q_1,\dots)\in (\mathbb{R}^{3N})^{\mathbb{N}}
\]
with conditional kernels
\[
K_{\Delta\tau}(q_{k+1}\mid q_0,\dots,q_k)
\propto
\exp\!\left(
-\frac{\Delta A(q_{k+1}\mid q_0,\dots,q_k)}{\hbar}
\right),
\]
where $\Delta A$ is the modal action increment.

These kernels are generically non--Markovian.

\subsection{Bohmian Limit}

In regimes where:
\begin{itemize}
  \item the action increment is sharply peaked,
  \item admissibility is stable,
  \item internal moduli decouple,
\end{itemize}
the kernels concentrate on the characteristic flow:
\[
K_{\Delta\tau}(\cdot\mid q_0,\dots,q_k)
\;\to\;
\delta\!\left(
q_{k+1} - \Phi^{(v)}_{\Delta\tau}(q_k)
\right),
\]
where $\Phi^{(v)}$ is the flow generated by
$v=j/\rho$.

\begin{theorem}[Pilot--Wave as Zero--Noise Limit]
Pilot--wave dynamics arise as the zero--noise,
peaked--kernel limit of the MTT indivisible stochastic
process within an admissible basin.
\end{theorem}

\begin{remark}
Outside this limit, deterministic Bohmian trajectories
are replaced by history--dependent stochastic evolution,
even though the underlying modal dynamics remain deterministic.
\end{remark}

%==================================================
\section{Summary of the MTT--to--Pilot--Wave Mapping}
%==================================================

\begin{center}
\begin{tabular}{|l|l|}
\hline
\textbf{Pilot--Wave Concept} & \textbf{MTT Origin} \\
\hline
Wavefunction $\psi$ &
Coherent--sector state under $P=\Pi\circ I$ \\[0.3em]
Born density $|\psi|^2$ &
Derived basin measure \\[0.3em]
Probability current $j$ &
Unitarity of $H_{\mathrm{obs}}$ \\[0.3em]
Guidance equation &
Characteristic flow of $j/\rho$ \\[0.3em]
Nonlocality &
Global action of coherent projector \\[0.3em]
Collapse &
Admissibility barrier crossing \\[0.3em]
Quantum equilibrium &
Consistency, not postulate \\[0.3em]
Global trajectories &
\emph{Forbidden} by noninvertibility \\[0.3em]
\hline
\end{tabular}
\end{center}

\bigskip
\noindent
\textbf{In Part IV (final)}, we will:
\begin{itemize}
  \item synthesize the results,
  \item contrast MTT, pilot--wave, and Everett interpretations,
  \item state precise conclusions and limits of applicability.
\end{itemize}
%==================================================
\section{Comparison with Interpretations of Quantum Mechanics}
%==================================================

We briefly situate the present derivation relative to other major interpretations
of quantum mechanics.

\subsection{Copenhagen-Type Interpretations}

Copenhagen interpretations treat the Schr\"odinger equation as fundamental
but supplement it with axiomatic measurement postulates.
Collapse is taken as primitive and non-unitary.

By contrast, in MTT:
\begin{itemize}
  \item the Schr\"odinger equation is derived rather than postulated,
  \item collapse is replaced by admissibility breakdown,
  \item no additional axioms beyond deterministic modal evolution are required.
\end{itemize}

\subsection{Everett / Many-Worlds}

Everettian interpretations retain universal unitary evolution
and deny physical collapse.
However, they require an additional interpretation of branching
and a nontrivial justification of the Born rule.

In MTT:
\begin{itemize}
  \item unitary evolution is only intra-basin,
  \item branching corresponds to admissibility basin structure,
  \item probabilities arise from basin measures, not branch counting.
\end{itemize}

Thus MTT agrees with Everett that collapse is not fundamental,
but rejects the claim of globally valid unitary evolution.

\subsection{Pilot--Wave Theory}

Pilot--wave theory restores determinism by introducing
configuration trajectories guided by the wavefunction.

The present paper shows that:
\begin{itemize}
  \item pilot--wave guidance equations are \emph{derived} in MTT,
  \item they require no additional ontological postulates,
  \item and they are valid only within admissible coherence basins.
\end{itemize}

MTT therefore subsumes pilot--wave theory as an effective description,
while explaining its scope limitations.

%==================================================
\section{Main Results}
%==================================================

We summarize the principal results of this paper.

\begin{enumerate}
  \item Modal Triplet Theory admits a deterministic modal flow
  on the full configuration space $X$.

  \item Coherent projection and observable pushforward
  induce an effective Schr\"odinger evolution
  on admissible coherence basins.

  \item The associated probability current defines
  a canonical configuration--space velocity field.

  \item The integral curves of this velocity field
  coincide with de Broglie--Bohm guidance equations.

  \item No hidden variables, stochastic postulates,
  or additional ontology are introduced.

  \item Pilot--wave dynamics are valid only intra--basin;
  they necessarily fail at admissibility barriers.

  \item Measurement corresponds to admissibility breakdown,
  not to physical collapse or branching.

  \item Outside the pilot--wave regime,
  MTT induces indivisible stochastic processes
  on observable histories.
\end{enumerate}

%==================================================
\section{Interpretive Statement}
%==================================================

The derivation presented here establishes the following:

\begin{quote}
\emph{Pilot--wave theory is not a fundamental alternative to quantum mechanics,
but an emergent characteristic description of Modal Triplet Theory
valid only within admissible coherence basins.}
\end{quote}

From the MTT perspective:
\begin{itemize}
  \item determinism is fundamental,
  \item probability is structural,
  \item nonlocal correlations arise from projection,
  \item and irreversibility is unavoidable at admissibility boundaries.
\end{itemize}

Pilot--wave mechanics correctly captures the semiclassical,
deterministic limit of this structure,
but cannot be globally extended without contradiction.

%==================================================
\section{Outlook}
%==================================================

Several directions follow naturally from this work.

\begin{itemize}
  \item \textbf{Relativistic extension:}
  Pilot--wave dynamics for relativistic fields arise
  from the same MTT framework without preferred foliations.

  \item \textbf{Quantum field theory:}
  Particle creation and annihilation correspond to
  basin transitions rather than trajectory endpoints.

  \item \textbf{Quantum gravity:}
  Horizon formation and evaporation are admissibility events,
  explaining information loss without violating determinism.

  \item \textbf{Cosmology:}
  Initial condition selection and probability measures
  are unified with quantum measurement under the same mechanism.
\end{itemize}

In this sense, Modal Triplet Theory provides a unifying
framework in which the successes of pilot--wave theory are retained,
its foundational problems are resolved,
and its limitations are rendered unavoidable and precise.
\begin{thebibliography}{99}

\bibitem{MTT-Admissibility}
P.~Nero,
\emph{Modal Triplet Theory: Admissibility, Encodings, and the Structure of Physical Description},
Zenodo preprint, January 2026.
\url{https://doi.org/10.5281/zenodo.18255621}

\bibitem{MTT-Foundation}
P.~Nero,
\emph{Modal Triplet Theory: Foundation},
Zenodo preprint, September 2025.
\url{https://doi.org/10.5281/zenodo.16949762}

\bibitem{MTT-FixedPoints}
P.~Nero,
\emph{Fixed Points I--VI: Complete Coherence Spine},
Zenodo preprints, August 2025.
\url{https://doi.org/10.5281/zenodo.16948748}

\bibitem{MTT-ProjectionAdmissibility}
P.~Nero,
\emph{The Projection--Admissibility Principle: Structural Constraints on Effective Physical Description},
Zenodo preprint, January 2026.
\url{https://doi.org/10.5281/zenodo.18255838}

\bibitem{MTT-Closure}
P.~Nero,
\emph{Closure and Inevitability in Modal Triplet Theory},
Zenodo preprint, January 2026.
\url{https://doi.org/10.5281/zenodo.18255510}

\bibitem{MTT-CoherenceCapacity}
P.~Nero,
\emph{Coherence Capacity as the Fundamental Resource of Effective Physics},
Zenodo preprint, January 2026.
\url{https://doi.org/10.5281/zenodo.18255905}

\bibitem{MTT-CoherenceDynamics}
P.~Nero,
\emph{Dynamics of Coherence Capacity: Transport, Concentration, and Exhaustion},
Zenodo preprint, January 2026.
\url{https://doi.org/10.5281/zenodo.18256048}

\bibitem{MTT-QM}
P.~Nero,
\emph{Modal Triplet Theory: From MTT to Quantum Mechanics},
Zenodo preprint, September 2025.
\url{https://doi.org/10.5281/zenodo.17074246}

\bibitem{MTT-QFT}
P.~Nero,
\emph{From MTT to Quantum Field Theory},
Zenodo preprint, 2025.
\url{https://doi.org/10.5281/zenodo.17068816}

\bibitem{MTT-GR}
P.~Nero,
\emph{Modal Triplet Theory: From MTT to General Relativity},
Zenodo preprint, October 2025.
\url{https://doi.org/10.5281/zenodo.16950597}

\bibitem{MTT-UVQG}
P.~Nero,
\emph{Modal Triplet Theory: From MTT to a UV-Finite, Unitary Quantum Gravity},
Zenodo preprint, 2025.
\url{https://doi.org/10.5281/zenodo.17077671}

\bibitem{MTT-Measurement}
P.~Nero,
\emph{Measurement as Disturbance and Stabilization in Modal Triplet Theory},
Zenodo preprint, 2025.
\url{https://doi.org/10.5281/zenodo.17177404}

\bibitem{MTT-Irreversibility}
P.~Nero,
\emph{Projection, Probability, and Irreversibility: Shadow Bridges Between Measurement, Black Holes, and Cosmology in Modal Triplet Theory},
Zenodo preprint, January 2026.
\url{https://doi.org/10.5281/zenodo.18256408}

\bibitem{MTT-Bell-Beables}
P.~Nero,
\emph{Modal Fixed Points, Bell’s Beables, and the Limits of Factorization},
Zenodo preprint, 2025.
\url{https://doi.org/10.5281/zenodo.17076300}

\bibitem{MTT-Bell-Temporal}
P.~Nero,
\emph{Temporal Bell Inequalities and Global Consistency in Modal Triplet Theory},
Zenodo preprint, August 2025.
\url{https://doi.org/10.5281/zenodo.18208884}

\bibitem{MTT-Stochastic}
P.~Nero,
\emph{From Modal Triplet Theory to Indivisible Stochastic Processes: A First-Principles, Fully Rigorous Derivation},
Zenodo preprint, January 2026.
\url{https://doi.org/10.5281/zenodo.18254862}

\end{thebibliography}
\end{document}
